\setlength{\parindent}{0pt}
% ------------------------------------------------------------
\chapter{Couplings and Probing Security}
\label{chap:couplings}
% ------------------------------------------------------------
\section{The Complexity of Deciding Secret-Independence}
\label{sec:hardness}

Definition~\ref{def:d-probing} phrases security as the condition that the law of a
probe stay fixed as the secrets vary, and Section~\ref{sec:hardness-preview}
already mentioned that this demand is actually a counting question. In this section we make
this claim precise. Underneath a single probe sit two problems, computing its law
and deciding whether varying the secret move it, and we determine each one's hardness. Computing the law
is a $\#\mathrm{P}$-complete count, and the decision problem is a
$\mathrm{C}_{=}\mathrm{P}$-complete comparison of such counts. This means
that verifying even one probe by counting is complete for classes that
sit above the entire polynomial hierarchy, which is the reason the thesis certifies
security constructively without explicit counting.

Throughout, fix a circuit over $\mathbb{F}_2$ with secret inputs $\mathcal{S}$,
random inputs $\mathcal{R}$, and public inputs $\mathcal{P}$, as usual, together with a probe
$\mathbf{w} = (w_1, \dots, w_k)$ of size $k \le d$. Under a secret assignment
$\mathbf{a} \in \{0,1\}^{\mathcal{S}}$ and a tape
$\rho \in \{0,1\}^{\mathcal{R}}$, the probe takes the joint value
$\mathbf{w}(\rho; \mathbf{a}) \in \{0,1\}^{k}$, with the publics held fixed.

  \subsection{From Independence to Distribution Counting}
  \label{sec:indep-to-counting}

  The law of a probe is assembled from integer counts. For an observation
  $o \in \{0,1\}^{k}$ write
  \[
    N_{\mathbf{a}}(o) \;=\; \big|\{\, \rho \in \{0,1\}^{\mathcal{R}} : \mathbf{w}(\rho; \mathbf{a}) = o \,\}\big| ,
  \]
  the number of tapes on which the probe reads $o$ under the secret $\mathbf{a}$.
  Since $\rho$ is uniform, the law is
  $\mathcal{L}_{\mathbf{a}}(\mathbf{w})(o) = N_{\mathbf{a}}(o) / 2^{|\mathcal{R}|}$, so the
  law is nothing but a histogram of these counts. Two laws coincide
  exactly when their counts do,
  \[
    \mathcal{L}_{\mathbf{a}}(\mathbf{w}) = \mathcal{L}_{\mathbf{b}}(\mathbf{w})
    \quad\Longleftrightarrow\quad
    N_{\mathbf{a}}(o) = N_{\mathbf{b}}(o) \ \text{ for every } o \in \{0,1\}^{k} .
  \]
  Secret-independence (Definition~\ref{def:d-probing}) is this equality quantified
  over all pairs of secrets. So what looks like a statement about distributions is
  actually a statement about integers, that certain counts of satisfying tapes agree.
  This splits into two questions of different character. One asks for the count itself;
  the other asks only whether two counts are equal, a decision. We treat them in turn.

  \begin{definition}[The counting and distinguishing problems]
  \label{def:count-distinguish}
  \textsc{Count} takes a circuit, a probe $\mathbf{w}$, a secret $\mathbf{a}$, and an
  observation $o$, and returns the integer $N_{\mathbf{a}}(o)$.
  \textsc{Distinguish} takes a circuit, a probe $\mathbf{w}$, and two secret assignments
  $\mathbf{a}, \mathbf{b}$, and accepts exactly when
  $\mathcal{L}_{\mathbf{a}}(\mathbf{w}) = \mathcal{L}_{\mathbf{b}}(\mathbf{w})$.
  \end{definition}

  \subsection{Counting the Law is \NoCaseChange{\#P}-Complete}
  \label{sec:sharpp}

  Recall that $\#\mathrm{P}$ is Valiant's class of functions counting the accepting
  computations of a nondeterministic polynomial-time machine, equivalently it is the
  enumeration of the satisfying assignments of a polynomial-size Boolean predicate
  \parencite{valiant1979complexity}.

  \begin{proposition}[Counting the law]
  \label{prop:count-sharpp}
  \textsc{Count} is $\#\mathrm{P}$-complete.
  \end{proposition}

  \begin{proof}
  \emph{Membership.} Fix $\mathbf{a}$, the publics, and $o$. The predicate
  $\rho \mapsto [\, \mathbf{w}(\rho; \mathbf{a}) = o \,]$ evaluates the circuit on the tape
  $\rho$ and compares $k$ output bits against $o$, all in time polynomial in the
  circuit size. Hence $N_{\mathbf{a}}(o)$ counts the tapes satisfying a
  polynomial-time predicate, so it is a $\#\mathrm{P}$ function.

  \emph{Hardness.} We reduce from $\#\textsc{SAT}$, the count of satisfying
  assignments of a Boolean formula, which is $\#\mathrm{P}$-complete
  \parencite{valiant1979complexity}. Given $\varphi(x_1, \dots, x_m)$, build the
  circuit that reads $x_1, \dots, x_m$ as random inputs, has no secrets, and
  computes $\varphi$ gate by gate, ending in a wire $w$ with $\wexpr{w} = \varphi$;
  the connectives of $\varphi$ are realised
  through $\oplus$ and $\odot$, which is possible because the two are universal over
  $\mathbb{F}_2$ (Section~\ref{sec:bool-circuits}). The one-wire probe $(w)$ then has
  $N(1) = |\{ x : \varphi(x) = 1 \}| = \#\varphi$, so any procedure for
  \textsc{Count} computes $\#\varphi$. The reduction is polynomial, hence
  \textsc{Count} is $\#\mathrm{P}$-hard.
  \end{proof}

  Notice that the hardness result uses no secret at all, so computing a probe's law is already
  intractable on the plainest circuits. That count is the object every exact method
  must, explicity or implicitly, realize.

  % The class wraps every \subsection title in titlecaps' \titlecap, which loops on
  % a math subscript; \oldsubsection bypasses that wrapper for this one heading.
  \oldsubsection{Deciding Independence is \texorpdfstring{\ensuremath{\mathrm{C}_{=}\mathrm{P}}}{C=P}-Complete}
  \label{sec:ceqp}

  \textsc{Distinguish} never asks for a count outright. It asks only whether two
  counts are equal, and that equality has a class of its own. Following Wagner,
  $\mathrm{C}_{=}\mathrm{P}$ is the class of languages of the form ``does a
  nondeterministic polynomial-time machine have exactly a given number of
  accepting paths'' \parencite{wagner1986complexity}. Another view comes
  from gap functions. Let $\mathrm{GapP}$ be the closure of $\#\mathrm{P}$ under
  subtraction, so the differences $f - g$ of two counting functions
  \parencite{fenner1994gap}. Then $\mathrm{C}_{=}\mathrm{P}$ is exactly the class of
  vanishing sets of gap functions,
  \[
    \big\{\, x \mid h(x) = 0 \,\big\}, \qquad h \in \mathrm{GapP} .
  \]
  Its canonical complete problem is \textsc{ExactSat}, deciding whether a formula
  has exactly a given number of satisfying assignments, equivalently one can also
  ask whether two formulas have the same number of satisfying assigmnets
  \parencite{wagner1986complexity}.

  \begin{theorem}[Deciding independence]
  \label{thm:distinguish-ceqp}
  \textsc{Distinguish} is $\mathrm{C}_{=}\mathrm{P}$-complete. Consequently deciding
  whether a probe is secret-independent is $\mathrm{C}_{=}\mathrm{P}$-hard.
  \end{theorem}

  \begin{proof}
  \emph{Membership.} Put
  \[
    f(C, \mathbf{w}, \mathbf{a}, \mathbf{b}) \;=\; \sum_{o \in \{0,1\}^{k}}
      \big(N_{\mathbf{a}}(o) - N_{\mathbf{b}}(o)\big)^2 ,
  \]
  Where $C$ is the circuit, each $N_{\mathbf{a}}$ is a $\#\mathrm{P}$ function,
  and $\mathrm{GapP}$ is closed under
  subtraction and multiplication \parencite{fenner1994gap}. The sum ranges over
  $2^{k} \le 2^{d}$ observations, which is a constant since the order $d$ is fixed, thus $f$ is
  a constant combination of gap functions and is itself in $\mathrm{GapP}$. A sum of
  squares vanishes exactly when every term does, hence
  \[
    \mathcal{L}_{\mathbf{a}}(\mathbf{w}) = \mathcal{L}_{\mathbf{b}}(\mathbf{w})
    \iff N_{\mathbf{a}}(o) = N_{\mathbf{b}}(o) \ \forall o
    \iff f(C, \mathbf{w}, \mathbf{a}, \mathbf{b}) = 0 ,
  \]
  and \textsc{Distinguish} is the vanishing set of the gap function $f$, so it lies in
  $\mathrm{C}_{=}\mathrm{P}$.

  \emph{Hardness.} We reduce from the equality form of \textsc{ExactSat}: given two
  formulas $\varphi_0, \varphi_1$ over $x_1, \dots, x_m$, decide whether
  $\#\varphi_0 = \#\varphi_1$, which is $\mathrm{C}_{=}\mathrm{P}$-complete
  \parencite{wagner1986complexity}. Build a circuit with a single secret bit $s$,
  the $x_i$ as randoms, and a final wire $w$ with
  \[
    \wexpr{w} \;=\; \varphi_0 \,\oplus\, \big(s \odot (\varphi_0 \oplus \varphi_1)\big) ,
  \]
  again realising the formulas over $\oplus$ and $\odot$. Then $w$ reads
  $\varphi_0$ when $s = 0$ and $\varphi_1$ when $s = 1$, so
  $N_{s=0}(1) = \#\varphi_0$ and $N_{s=1}(1) = \#\varphi_1$. The one-bit probe
  $(w)$ satisfies $\mathcal{L}_{0} = \mathcal{L}_{1}$ if and only if
  $\#\varphi_0 = \#\varphi_1$, so \textsc{Distinguish} on this circuit with
  $\mathbf{a} = 0$, $\mathbf{b} = 1$ decides the equality. The construction is
  polynomial, so \textsc{Distinguish} is $\mathrm{C}_{=}\mathrm{P}$-hard. One secret
  bit and one probed wire already suffice, so no feature of the masking is what
  makes the problem hard.
  \end{proof}

  \subsection{What the Placements Mean for Verification}
  \label{sec:hardness-consequences}

  The two results are closely related. The law's entries are $\#\mathrm{P}$
  counts, and the security question is whether two of them coincide, which is the
  defining shape of the $\mathrm{C}_{=}\mathrm{P}$ class. The classes fall in the expected
  places. Every $\mathrm{C}_{=}\mathrm{P}$ query is settled by two calls to a
  counting oracle, followed by a subtraction, and a test against zero, so
  $\mathrm{C}_{=}\mathrm{P} \subseteq \mathrm{P}^{\#\mathrm{P}}$. Which means
  that the decision never costs more than the counting.

  Another observation is the following: a formula is unsatisfiable exactly when its count is
  zero (as an instance of exact counting) so $\mathrm{coNP} \subseteq
  \mathrm{C}_{=}\mathrm{P}$. Hence \textsc{Distinguish} is already
  $\mathrm{coNP}$-hard, and no polynomial-time exact decision exists unless
  $\mathrm{P} = \mathrm{NP}$. Further showing us that deciding
  one probe exactly is very much intractable.

  The full verification task only adds on this. $d$-probing security asks \textsc{Distinguish}
  to answer ``equal'' for every one of the $\sum_{i \le d} \binom{|\mathcal{W}|}{i}$
  probes and every pair of secrets. A universal quantifier over exponentially many
  instances of a $\mathrm{C}_{=}\mathrm{P}$-complete question.

  So now we should be fairly certain that computing or
  equating these counts is out of reach, so we do not compute them. We certify security
  constructively instead, by exhibiting for a probe a measure-preserving relabelling
  of the randomness under which its secrets visibly cancel. Such a witness proves independence without ever
  counting, though at the price of completeness, since a probe can be independent for
  reasons no relabelling of our kind captures.

\section{Couplings of Probabilistic Programs}
\label{sec:couplings-prob}

The hardness of Section~\ref{sec:hardness} rules out counting, so we need another
way of certifying that two laws agree. The workaround idea comes from the fact that 
our two laws are not two unrelated distributions. They come from the same uniform
random source.

Fix a probe $\mathbf{w}$ and two secret assignments $\mathbf{a}$ and $\mathbf{b}$.
Both $\mathcal{L}_{\mathbf{a}}(\mathbf{w})$ and
$\mathcal{L}_{\mathbf{b}}(\mathbf{w})$ are produced the same way, by drawing one
uniform tape $\rho \in \{0,1\}^{\mathcal{R}}$ and reading the probe; the only
difference between them is which secret is plugged in. So rather than computing the
two distributions and comparing them, we can try to match up the tapes that produce
them.

Suppose we can pair each tape $\rho$ with a tape $\rho'$ so that the run under
$\mathbf{a}$ on $\rho$ and the run under $\mathbf{b}$ on $\rho'$ read exactly the
same values,
\[
  \mathbf{w}(\rho; \mathbf{a}) \;=\; \mathbf{w}(\rho'; \mathbf{b}) ,
\]
and suppose the pairing $\rho \mapsto \rho'$ is a bijection of the tape space. Then
the two laws are equal. Because, the bijection matches
the tapes producing an observation $o$ under $\mathbf{a}$ one-for-one with the tapes
producing $o$ under $\mathbf{b}$, so the counts $N_{\mathbf{a}}(o)$ and
$N_{\mathbf{b}}(o)$ of Section~\ref{sec:indep-to-counting} agree without either
being computed explicitly. A relabelling of the randomness has replaced counting, which is
exactly the trade we were looking for. Such a bijection is what we call a coupling.

The name is borrowed from probability theory, where a \emph{coupling} of two
distributions $\mu$ and $\nu$ on a set $X$ is a joint distribution on
$X \times X$ whose marginals are $\mu$ and $\nu$, and where a coupling supported on
the diagonal forces $\mu = \nu$. This general notion is what underlies the coupling-based
program logics of Barthe et al.\ \parencite{barthe2017coupling}.

  \subsection{Coupling Two Runs Under Different Secrets}
  \label{sec:coupling-two-runs}

  \begin{definition}[Coupling function]
  \label{def:coupling-function}
  Let $\mathbf{w}$ be a probe and $\mathbf{a}, \mathbf{b} \in \{0,1\}^{\mathcal{S}}$ two
  secret assignments. A \emph{coupling function} from $\mathbf{a}$ to $\mathbf{b}$
  for $\mathbf{w}$ is a bijection $\varphi : \{0,1\}^{\mathcal{R}} \to \{0,1\}^{\mathcal{R}}$
  with
  \[
    \mathbf{w}(\rho; \mathbf{a}) \;=\; \mathbf{w}\big(\varphi(\rho); \mathbf{b}\big)
    \qquad \text{for every } \rho \in \{0,1\}^{\mathcal{R}} .
  \]
  \end{definition}

  A bijection of the finite set $\{0,1\}^{\mathcal{R}}$ automatically preserves its
  uniform law, so ``measure-preserving'' actually just means bijectivity here, as
  Section~\ref{sec:prob-independence} already noted. What
  the identity demands is the key point. It says $\varphi$ relabels the tapes so
  that a run under $\mathbf{a}$ is reproduced, observation for observation, by a
  run under $\mathbf{b}$. Such a $\varphi$ is precisely a diagonal coupling of the
  two laws, realised by the map $\rho \mapsto \big(\mathbf{w}(\rho; \mathbf{a}),\,
  \mathbf{w}(\varphi(\rho); \mathbf{b})\big)$, whose image lies on the diagonal.

  Collecting one function per ordered pair of secrets gives a family
  $\Phi = \{\, \Phi_{\mathbf{a},\mathbf{b}} \,\}_{\mathbf{a},\mathbf{b}}$ of coupling
  functions. We may always take $\Phi_{\mathbf{a},\mathbf{a}} = \mathrm{id}$,
  $\Phi_{\mathbf{b},\mathbf{a}} = \Phi_{\mathbf{a},\mathbf{b}}^{-1}$, and
  $\Phi_{\mathbf{b},\mathbf{c}} \circ \Phi_{\mathbf{a},\mathbf{b}} =
  \Phi_{\mathbf{a},\mathbf{c}}$, so the family is a groupoid on the secret
  assignments and one row $\{\Phi_{\mathbf{a}_0,\,\cdot}\}$ generates the rest.

\section{Probing Security via Couplings}
\label{sec:probing-via-couplings}

The definition is engineered so that its existence and secret-independence are one
and the same statement.

\begin{theorem}[Security is a coupling]
\label{thm:coupling-equiv}
A probe $\mathbf{w}$ is secret-independent if and only if for every ordered pair of secret
assignments $\mathbf{a}, \mathbf{b}$ there is a coupling function from $\mathbf{a}$
to $\mathbf{b}$ for $\mathbf{w}$.
\end{theorem}

\begin{proof}
Suppose $\varphi$ is a coupling function from
$\mathbf{a}$ to $\mathbf{b}$. For any observation $o$, the identity
$\mathbf{w}(\rho; \mathbf{a}) = \mathbf{w}(\varphi(\rho); \mathbf{b})$ and the bijectivity of
$\varphi$ give
\[
  N_{\mathbf{a}}(o) = \big|\{\rho : \mathbf{w}(\rho; \mathbf{a}) = o\}\big|
  = \big|\{\rho : \mathbf{w}(\varphi(\rho); \mathbf{b}) = o\}\big|
  = \big|\{\rho' : \mathbf{w}(\rho'; \mathbf{b}) = o\}\big| = N_{\mathbf{b}}(o) ,
\]
reindexing by $\rho' = \varphi(\rho)$. Hence
$\mathcal{L}_{\mathbf{a}}(\mathbf{w}) = \mathcal{L}_{\mathbf{b}}(\mathbf{w})$, and if this holds for
every pair then $\mathbf{w}$ is secret-independent by definition.

For the other way, suppose
$\mathcal{L}_{\mathbf{a}}(\mathbf{w}) = \mathcal{L}_{\mathbf{b}}(\mathbf{w})$, i.e.\
$N_{\mathbf{a}}(o) = N_{\mathbf{b}}(o)$ for every $o$. The sets
$F_{\mathbf{a}}(o) = \{\rho \mid \mathbf{w}(\rho; \mathbf{a}) = o\}$ partition
$\{0,1\}^{\mathcal{R}}$ as $o$ ranges over the observations, and likewise the
$F_{\mathbf{b}}(o)$. Corresponding blocks have equal size, so we may pick a
bijection $F_{\mathbf{a}}(o) \to F_{\mathbf{b}}(o)$ for each $o$ and take their
union $\varphi$. It is a bijection of $\{0,1\}^{\mathcal{R}}$, and by construction
it carries each $\rho \in F_{\mathbf{a}}(o)$ to some
$\rho' \in F_{\mathbf{b}}(o)$, so $\mathbf{w}(\varphi(\rho); \mathbf{b}) = o =
\mathbf{w}(\rho; \mathbf{a})$. So, $\varphi$ is a coupling function from $\mathbf{a}$ to
$\mathbf{b}$.
\end{proof}

  \subsection{\NoCaseChange{$d$}-Probing Security as a Family of Couplings}
  \label{sec:dprobing-couplings}

  Applying Theorem~\ref{thm:coupling-equiv} to every probe lifts the equivalence
  to the whole circuit. It is $d$-probing secure exactly when every probe of size
  at most $d$ admits a family of coupling functions, one per ordered pair of
  secrets. Two runs with differing secrets give any $d$ wires the same law
  precisely because a coupling can relabel the randomness of one into the
  randomness of the other, wire for wire.

  The construction of such a coupling, assuming that the probe does not leak, turns
  the assumption into a coupling by matching fibres block by block, and to do so it needs
  the counts $N_{\mathbf{a}}(o)$, which are the very quantities Section~\ref{sec:hardness}
  proved to be intractable. So the equivalence, is an existence
  statement and not really a recipe. It guarantees a coupling without telling us how to
  find one nor how to check it efficiently. Supplying such an explicit
  witness is the task of the construction we explain below.

\section{Substitutions as a Coupling Construction}
\label{sec:subst-coupling}

The randomness substitutions of Section~\ref{sec:randomness-substitution} can become the
building blocks of an explicit coupling. Recall that a substitution
$\sigma = (r \mapsto e \oplus r)$ induces a bijection
$\widehat{\sigma}$ of $\{0,1\}^{\mathcal{R}}$, with
$(w\sigma)(\rho) = w(\widehat{\sigma}(\rho))$ for every wire
(Lemma~\ref{lem:substitution}). Each such $\widehat{\sigma}$ is a single relabelling
of the tape, and relabellings of this shape
compose. A finite chain of substitutions therefore composes into a bijection of the
randomness, and it is this composite that we will exhibit as the witness of
Theorem~\ref{thm:coupling-equiv}.

  \subsection{The Derivation of an Explicit Coupling Witness}
  \label{sec:derivation-witness}

  \begin{definition}[Derivation]
  \label{def:derivation}
  A \emph{derivation} for a probe $\mathbf{w}$ is a finite sequence of
  substitutions
  \[
    \sigma_1 = (r_1 \mapsto e_1 \oplus r_1), \ \dots, \
    \sigma_K = (r_K \mapsto e_K \oplus r_K),
    \qquad r_i \notin \mathrm{vars}(e_i) ,
  \]
  applied in order. It \emph{ends secret-free} when the resulting expression
  \[
    \overline{\mathbf{w}} \;=\; \mathbf{w}\sigma_1 \sigma_2 \cdots \sigma_K
  \]
  is secret-free in the sense of Notation~\ref{not:fv-vars}, i.e.\ no secret
  variable occurs in it.
  \end{definition}

  Before we construct our coupling from the derivation we mention two small details.
  First, Lemma~\ref{lem:substitution} was stated with the secrets held fixed, however now
  we let them vary, and we must record how the pieces depend on them. Second, the
  context $e_i$ of a substitution may itself mention secret variables, so the shift
  it applies to the $r_i$-coordinate is not one map but one map per secret
  assignment. We write $\widehat{\sigma_i}^{\,\mathbf{a}}$ for the bijection
  induced by $\sigma_i$ under the secret assignment $\mathbf{a}$, the one that
  shifts the $r_i$-coordinate by $e_i(\rho; \mathbf{a})$.

  Composing these under one secret assignment $\mathbf{a}$, let
  \[
    \varphi_{\mathbf{a}} \;=\; \widehat{\sigma_1}^{\,\mathbf{a}} \circ
      \widehat{\sigma_2}^{\,\mathbf{a}} \circ \cdots \circ
      \widehat{\sigma_K}^{\,\mathbf{a}} ,
  \]
  a bijection of $\{0,1\}^{\mathcal{R}}$, being a composition of bijections. The
  derivation is thus a purely syntactic list of rewrite steps,
  which yields an actual tape bijection only once a secret assignment is supplied.
  The theorem below says that these bijections fit together into a coupling.

  \begin{theorem}[Derivations yield couplings]
  \label{thm:derivation-coupling}
  Let $\sigma_1, \dots, \sigma_K$ be a derivation for a probe $\mathbf{w}$ that
  ends secret-free, and let $\varphi_{\mathbf{a}}$ be as above. Then for every
  ordered pair of secret assignments
  \[
    \Phi_{\mathbf{a},\mathbf{b}} \;=\; \varphi_{\mathbf{b}} \circ
      \varphi_{\mathbf{a}}^{-1}
  \]
  is a coupling function from $\mathbf{a}$ to $\mathbf{b}$ for $\mathbf{w}$ in the
  sense of Definition~\ref{def:coupling-function}. Consequently $\mathbf{w}$ is
  secret-independent.
  \end{theorem}

  \begin{proof}
  Fix a secret assignment $\mathbf{a}$ for the moment. Lemma~\ref{lem:substitution}
  says that rewriting the expression by one substitution is the same as moving the
  tape by the bijection it induces. Doing this once per step and composing gives
  \[
    \mathbf{w}\big(\varphi_{\mathbf{a}}(\rho);\, \mathbf{a}\big)
      \;=\; \overline{\mathbf{w}}(\rho;\, \mathbf{a})
    \qquad \text{for every tape } \rho .
  \]
  This can be shown formally with an induction on $K$.
  % and the order is worth noting: the step
  % applied last to the expression is the one applied first to the tape, which is why
  % $\varphi_{\mathbf{a}}$ composes the shears in that order.

  Since $\overline{\mathbf{w}}$ is secret-free, evaluating
  it never reads the secret assignment, so we may write $\overline{\mathbf{w}}(\rho)$
  and read the identity above at $\mathbf{a}$ and at $\mathbf{b}$, as follows:
  \[
    \mathbf{w}\big(\varphi_{\mathbf{a}}(\rho);\, \mathbf{a}\big)
      \;=\; \overline{\mathbf{w}}(\rho) \;=\;
    \mathbf{w}\big(\varphi_{\mathbf{b}}(\rho);\, \mathbf{b}\big) .
  \]
  Both runs are relabelled views of the same secret-free expression, and it only
  remains to link them to each other. As $\varphi_{\mathbf{a}}$ is a bijection we may
  put $\varphi_{\mathbf{a}}^{-1}(\rho)$ in place of $\rho$, which leaves
  $\mathbf{w}(\rho; \mathbf{a})$ on the left and gives
  \[
    \mathbf{w}(\rho; \mathbf{a}) \;=\;
    \mathbf{w}\big(\Phi_{\mathbf{a},\mathbf{b}}(\rho);\, \mathbf{b}\big)
    \qquad \text{for every } \rho ,
  \]
  with $\Phi_{\mathbf{a},\mathbf{b}} = \varphi_{\mathbf{b}} \circ
  \varphi_{\mathbf{a}}^{-1}$ being a bijection, as a composition of two. This is
  exactly Definition~\ref{def:coupling-function}. The pair $\mathbf{a}, \mathbf{b}$
  was arbitrary, so Theorem~\ref{thm:coupling-equiv} applies and $\mathbf{w}$ is
  secret-independent.
  \end{proof}

  The family $\{\Phi_{\mathbf{a},\mathbf{b}}\}$ built this way satisfies the
  groupoid identities of Section~\ref{sec:coupling-two-runs} for free, since
  $\Phi_{\mathbf{a},\mathbf{a}} = \mathrm{id}$,
  $\Phi_{\mathbf{b},\mathbf{a}} = \Phi_{\mathbf{a},\mathbf{b}}^{-1}$ and
  $\Phi_{\mathbf{b},\mathbf{c}} \circ \Phi_{\mathbf{a},\mathbf{b}} =
  \Phi_{\mathbf{a},\mathbf{c}}$ all follow from the shape of
  $\varphi_{\mathbf{b}} \circ \varphi_{\mathbf{a}}^{-1}$.

  This is the constructive certificate that we wished for.
  It is explicit, constructed straight off the recorded
  substitutions, and it is checkable without ever counting a fibre.
  % What it is not
  % is \emph{complete}, since a probe can be secret-independent while no chain of
  % substitutions reaches a secret-free form, in which case the search returns no
  % witness, though one might exist in the abstract sense of
  % Theorem~\ref{thm:coupling-equiv}.
  How such a derivation is searched for, recorded,
  and then reused to certify many wires at once is the subject of
  the next chapter (Chapter~\ref{chap:design}).
